\documentclass[final]{article}
\usepackage[paperwidth=84.1cm, paperheight=118.9cm, margin=2.5cm]{geometry}
\usepackage{multicol}
\usepackage{xcolor}
\usepackage{amsmath,amssymb}
\usepackage{tcolorbox}
\usepackage{graphicx}
\usepackage{array}

\definecolor{primary}{RGB}{255, 107, 53}
\definecolor{secondary}{RGB}{15, 23, 42}
\definecolor{accent}{RGB}{6, 182, 212}
\definecolor{boxbg}{RGB}{248, 250, 252}

\tcbset{
  colback=boxbg,
  colframe=primary,
  coltitle=white,
  fonttitle=\Huge\bfseries,
  boxrule=3mm,
  arc=6mm,
  boxsep=1.5cm,
  left=1cm, right=1cm, top=1cm, bottom=1cm
}

\begin{document}
\pagestyle{empty}

% HEADER GIGANTE DEL PÓSTER
\begin{tcolorbox}[colback=secondary, colframe=primary, boxrule=4mm, arc=8mm]
  \centering
  \vspace{0.8cm}
  {\fontsize{64}{76}\selectfont \bfseries \color{white} TÍTULO PRINCIPAL DEL PÓSTER CIENTÍFICO}\\[1.5cm]
  {\fontsize{38}{46}\selectfont \color{primary} \textbf{Dr. Nombre del Investigador}$^1$, \textbf{MSc. Co-Autor Uno}$^2$, \textbf{Lic. Co-Autor Dos}$^1$}\\[1cm]
  {\fontsize{28}{36}\selectfont \color{gray!40} $^1$Departamento de Ciencias e Ingeniería, Universidad Nacional \quad $\bullet$ \quad $^2$Laboratorio de Investigación Avanzada}\\[0.6cm]
  {\fontsize{24}{30}\selectfont \color{accent} Contacto: \texttt{investigador@universidad.edu} \quad $\bullet$ \quad Congreso Internacional de Ciencia y Tecnología 2026}
  \vspace{0.8cm}
\end{tcolorbox}

\vspace{1.5cm}

\setlength{\columnsep}{2.5cm}
\begin{multicols}{3}

  % COLUMNA 1
  \begin{tcolorbox}[title=1. Introducción y Motivación]
    {\LARGE
    El procesamiento eficiente de grandes volúmenes de datos requiere arquitecturas optimizadas que minimicen los tiempos de latencia y el consumo energético.
    
    \vspace{0.8cm}
    \textbf{Objetivos Principales:}
    \begin{itemize}
      \setlength{\itemsep}{0.6cm}
      \item Desarrollar un modelo analítico escalable.
      \item Validar la convergencia numérica en tiempo real.
      \item Evaluar el rendimiento frente a estándares actuales.
    \end{itemize}
    }
  \end{tcolorbox}

  \vspace{1.2cm}

  \begin{tcolorbox}[title=2. Marco Teórico y Formulación]
    {\LARGE
    La dinámica fundamental del sistema se expresa a través de la ecuación diferencial generalizada de conservación:
    
    \begin{equation*}
      \frac{\partial \rho}{\partial t} + \nabla \cdot (\rho \mathbf{u}) = \sum_{k=1}^{N} \mathcal{S}_k(\mathbf{x}, t)
    \end{equation*}

    Donde $\rho$ representa la densidad de probabilidad del estado y $\mathbf{u}$ el campo vectorial de transporte acoplado.
    }
  \end{tcolorbox}

  % COLUMNA 2
  \begin{tcolorbox}[title=3. Metodología Experimental]
    {\LARGE
    Se implementó un esquema de muestreo discreto con validación cruzada cuádruple:
    
    \vspace{0.8cm}
    \begin{enumerate}
      \setlength{\itemsep}{0.6cm}
      \item Adquisición de señales a frecuencia de muestreo $f_s = 48\text{ kHz}$.
      \item Filtrado adaptativo mediante transformada wavelet ortogonal.
      \item Optimización convexa del vector de parámetros $\boldsymbol{\theta}^*$.
    \end{enumerate}
    }
  \end{tcolorbox}

  \vspace{1.2cm}

  \begin{tcolorbox}[title=4. Resultados y Métricas]
    {\LARGE
    Los ensayos experimentales confirman una reducción sustancial del margen de error:
    
    \vspace{1cm}
    \centering
    \renewcommand{\arraystretch}{1.6}
    \begin{tabular}{|l|c|c|c|}
      \hline
      \textbf{Algoritmo} & \textbf{Precisión} & \textbf{Tiempo (s)} & \textbf{Speedup} \\
      \hline
      Modelo Base & 88.4\% & 14.8 & $1.0\times$ \\
      Propuesta A & 94.1\% & 6.2 & $2.4\times$ \\
      \textbf{NaTex Model} & \textbf{99.2\%} & \textbf{1.9} & $\mathbf{7.8\times}$ \\
      \hline
    \end{tabular}
    \vspace{0.8cm}
    }
  \end{tcolorbox}

  % COLUMNA 3
  \begin{tcolorbox}[title=5. Discusión y Análisis]
    {\LARGE
    La tasa de convergencia lograda evidencia que la formulación matricial reduce la complejidad temporal de $\mathcal{O}(N^3)$ a $\mathcal{O}(N \log N)$.
    
    \vspace{0.8cm}
    Las perturbaciones térmicas y de ruido Gaussiano se mitigaron con una relación señal-ruido ($\text{SNR}$) superior a $38\text{ dB}$.
    }
  \end{tcolorbox}

  \vspace{1.2cm}

  \begin{tcolorbox}[title=6. Conclusiones y Referencias]
    {\LARGE
    \textbf{Conclusiones:}
    \begin{itemize}
      \setlength{\itemsep}{0.5cm}
      \item Se validó la reproducibilidad del modelo al 100\%.
      \item El método es apto para despliegue industrial.
    \end{itemize}

    \vspace{1cm}
    \textbf{Referencias Clave:}
    \begin{enumerate}
      \setlength{\itemsep}{0.4cm}
      \item Shannon, C. E. (1948). \textit{A Mathematical Theory of Communication}.
      \item Goodfellow, I. et al. (2016). \textit{Deep Learning}, MIT Press.
    \end{enumerate}
    }
  \end{tcolorbox}

\end{multicols}

\end{document}
